% =============================================================================
% BAB I PENDAHULUAN -- teks panduan templat Skripsi Sain-Tek-Kes 2026
% Ganti seluruh teks panduan dengan isi skripsi.
% =============================================================================

\chapter{PENDAHULUAN}

\begin{comment}

Bab pendahuluan memuat latar belakang atau justifikasi dipilihnya topik karya
ilmiah tugas akhir, perumusan atau pendekatan penyelesaian masalah, tujuan,
manfaat, dan ruang lingkup. Di dalam pendahuluan dijelaskan perumusan atau
pendekatan penyelesaian masalah dan alasan pemilihan metode yang digunakan.
Merujuk pada proses perumusan masalah tugas akhir, bagian Kerangka Pikir dan
Hipotesis dapat ditulis di sini, tidak ditulis dalam bab tersendiri. Untuk
membantu mengikuti alur pikir secara skematis, dapat juga dibuat bagan alir
kerangka proses dan rumusan masalah serta pencapaian tujuan tugas akhir.
\end{comment}

\section{Latar Belakang}

% Dicek 2026-11-21 di justifikasi asesmen IUCN (species/19488/126673248/pdf):
% "a population reduction of 50-79% was inferred over the last three
% generations (120 years). Therefore, the species is assessed as Endangered
% A2bd." Indo-Pasifik: penurunan >65%; Atlantik: >=30%.

Populasi hiu paus (\textit{Rhincodon typus}) global diperkirakan turun 50--79\%
dalam tiga generasi terakhir (120 tahun), sehingga spesies ini berstatus
\textit{Endangered} dalam Daftar Merah IUCN \citep{Pierce2025}. Di
Indo-Pasifik, wilayah tempat perairan Indonesia berada, penurunan populasi
lebih dari 65\%, sedangkan di Atlantik sedikitnya 30\% \citep{Pierce2025}. Hiu
paus dapat mencapai panjang total 18--20 m, mampu bergerak jauh di samudra
tropis dan subtropis hangat, dan terutama terancam oleh tangkapan sampingan
jaring perikanan serta tabrakan dengan kapal \citep{Pierce2025}. Perlindungan di lokasi agregasi belum cukup karena
pergerakan hiu paus melampaui kawasan yang dikelola dan banyak habitat berada di
luar area konservasi \citep{Reynolds2017,Putra2026}. Oleh karena itu,
pengelolaan membutuhkan peta habitat yang sesuai. Peta seperti ini dapat
disusun dengan model kesesuaian habitat \citep{Putra2026}.

Model kesesuaian habitat hiu paus telah dibangun di berbagai wilayah, dengan
data dari telemetri, \textit{logbook} perikanan, atau penampakan yang
dilaporkan masyarakat setempat \citep{Reynolds2017,Sequeira2014,Milles2025}. Di
Australia Barat, \citet{Reynolds2017} memodelkan habitat hiu paus dari data
telemetri hiu yang ditandai di Ningaloo selama 2010--2016. Model yang
digunakan adalah model aditif campuran (GAMM) dengan peubah kedalaman, suhu
permukaan laut, dan klorofil-a. Model tersebut memprediksi habitat yang sesuai
hingga Laut Timor dan perairan Indonesia.
% DRAF ASISTEN 2026-11-21 (diminta penulis): mulai kalimat Miller sampai akhir
% paragraf. Sunting ke gaya sendiri.
Di perairan timur Australia hingga Pasifik barat daya, \citet{Miller2026}
membangun model ensembel dari data telemetri hiu paus muda. Kinerja model
tersebut tinggi pada data latih (AUC 0,98), tetapi turun menjadi 0,65--0,69
ketika diuji dengan validasi silang blok spasial. Pada skala samudra, data
pengamatan dari kapal pukat cincin tuna telah digunakan untuk memodelkan
sebaran hiu paus di Samudra Atlantik \citep{Baez2020} serta di Samudra
Atlantik, Hindia, dan Pasifik sekaligus \citep{Sequeira2014}. Kajian global
tersebut menemukan bahwa pengaruh peubah lingkungan terhadap kemunculan hiu
paus memiliki sedikit perbedaan antarsamudra \citep{Sequeira2014}. Data penampakan dari
sumber data terbuka dan laporan masyarakat juga telah digunakan, antara lain di
Teluk California \citep{PetatanRamirez2020} dan Atlantik barat laut
\citep{Milles2025}. Kajian-kajian tersebut dibangun di wilayah dengan
ketersediaan data yang tinggi, yaitu wilayah dengan telemetri jangka panjang,
catatan perikanan selama puluhan tahun, atau penampakan yang intensif.
Model Australia tersebut memprediksi habitat yang sesuai hingga perairan
Indonesia, tetapi hanya divalidasi dengan hiu yang ditandai di Ningaloo, tanpa
data kemunculan independen dari Indonesia \citep{Reynolds2017}.

% DRAF P3 (Gemini Neon, 2026-11-21, dirapikan asisten Claude Opus 5.5): sunting ke gaya sendiri.
Berbeda dari wilayah-wilayah tersebut, kajian hiu paus di Asia Tenggara umumnya
berskala lokal, berfokus pada satu lokasi agregasi, dan tidak semuanya berupa
model habitat. Di Filipina, \citet{Araujo2019} mengidentifikasi individu hiu
paus di Teluk Honda, Palawan, melalui identifikasi citra dan penanda satelit
tanpa membangun model habitat. Kajian tersebut menunjukkan bahwa hiu paus
melintasi batas negara. Seekor hiu yang mula-mula tercatat di Kalimantan Timur,
Indonesia, terlihat kembali di Teluk Honda sekitar 3,5 tahun kemudian. Hiu
bertanda lain mencapai Sabah, Malaysia, lalu kembali \citep{Araujo2019}. Di
Indonesia, kajian di Teluk Cenderawasih menerapkan GAM pada data 2012--2015
\citep{Ranintyari2018} dan 2019--2020 \citep{Manuhutu2021}. Kajian lain di
teluk yang sama meneliti populasi hiu paus serta suhu dan salinitas perairan
\citep{Adriani2023}. Di Teluk Saleh, kajian menggunakan regresi linear
sederhana dari 36 penampakan \citep{Herdianti2025}. Kajian-kajian lokal
tersebut menemukan bahwa suhu permukaan laut dan klorofil-a berkaitan dengan
kemunculan hiu paus \citep{Ranintyari2018,Manuhutu2021,Herdianti2025}. Arus
permukaan juga berpengaruh di Teluk Cenderawasih
\citep{Ranintyari2018,Manuhutu2021}. Pada skala kepulauan, \citet{Putra2026}
menggunakan telemetri satelit dari 70 hiu paus yang seluruhnya ditandai di
empat lokasi agregasi (Teluk Cenderawasih, Kaimana, Teluk Saleh, dan Teluk
Tomini). Data tersebut digunakan untuk membangun 73 model MaxEnt dengan
rata-rata AUC 0,82. Sementara itu, \citet{Setyawan2025} mencatat 1.118
penampakan selama 13 tahun di Bentang Laut Kepala Burung. Hampir seluruh
penampakan tersebut berada di sekitar bagan, dan 97,8\% di antaranya hanya
berasal dari dua wilayah, yaitu Teluk Cenderawasih dan Kaimana. Populasinya
didominasi jantan muda dengan rasio 8,8:1 \citep{Setyawan2025}. Dengan
demikian, data kemunculan hiu paus di Indonesia memiliki ketersediaan yang
rendah, mengelompok di beberapa lokasi agregasi, dan bias terhadap jantan muda.
Akibatnya, model yang hanya dilatih dengan data Indonesia berisiko tidak
mewakili seluruh perairan Indonesia.

% DRAF P4 (Gemini Neon, 2026-11-21, dirapikan asisten Claude Opus 5.5): sunting ke gaya sendiri.
Salah satu cara untuk mengatasi keterbatasan data Indonesia adalah memanfaatkan
data dari wilayah dengan ketersediaan data yang tinggi, seperti Australia.
Pilihan ini beralasan karena model Australia sudah memprediksi habitat yang
sesuai di perairan Indonesia \citep{Reynolds2017}. Pendekatan serupa pernah
diterapkan oleh \citet{Bridges2023}, yang mentransfer model MaxEnt karang air
dingin dari Atlantik Utara yang kaya data ke Atlantik Selatan yang kekurangan data
dan memiliki kondisi lingkungan serupa. Model hasil transfer tersebut mencapai
AUC 0,70 pada data validasi independen. Namun, \citet{Wilson2026} menemukan
bahwa tidak satu pun dari 24 konfigurasi model berbasis telemetri sepuluh
spesies burung laut dan mamalia laut di Samudra Selatan mencapai kemampuan
transfer spasial yang tinggi secara konsisten. Mereka menjelaskan bahwa model
yang ditransfer lintas ruang umumnya berkinerja buruk karena terjadi
ekstrapolasi lingkungan dan relung lingkungan dapat berbeda antarpopulasi. Di
Indonesia, \citet{Sahri2023} memproyeksikan model \textit{Random Forest} paus
dan lumba-lumba dari empat wilayah kalibrasi ke seluruh Zona Ekonomi Eksklusif
Indonesia. Proyeksi untuk paus sperma menempatkan kesesuaian tinggi di dekat
pesisir pulau, padahal habitatnya berada di laut lepas. Para penulis tersebut
mengaitkan hal ini dengan bias pengambilan sampel ke arah pesisir pada
penampakan oportunistis. Secara umum, model yang berkinerja serupa di wilayah
kalibrasi dapat berbeda jauh ketika ditransfer ke kondisi lingkungan baru
\citep{Elith2010}. Selain perbedaan lingkungan, kajian global pada hiu paus
juga menunjukkan bahwa pengaruh peubah lingkungan terhadap kemunculan sedikit
berbeda antarsamudra \citep{Sequeira2014}. Oleh karena itu, data Australia
tidak dapat digunakan begitu saja. Pemodelan habitat membutuhkan cara untuk
menangani perbedaan antara wilayah sumber (Australia) dan wilayah target
(Indonesia).

% DRAF P5 (Gemini Neon, 2026-11-21, dirapikan asisten): sunting ke gaya sendiri.
Selain transfer langsung, beberapa pendekatan dapat memanfaatkan data dari
wilayah sumber, dan setiap pendekatan menangani masalah yang berbeda.
\textit{Subsampling} lingkungan, seperti pengelompokan lingkungan dan
penipisan jarak lingkungan, menyaring catatan kemunculan pada ruang lingkungan
untuk mengoreksi bias pengambilan sampel \citep{Pili2025}. Pada uji spesies
virtual, pendekatan ini menghasilkan model paling tangguh untuk mentransfer
prediksi ke lingkungan baru \citep{Pili2025}. Namun, pendekatan ini tidak
menggunakan informasi tentang wilayah target, sehingga penelitian ini
memakainya sebagai pembanding. Pendekatan lain adalah penggabungan data
\textit{presence-absence} dan \textit{presence-only}. Dalam simulasi,
penggabungan tersebut meningkatkan transfer model kesesuaian habitat ke lokasi
yang minim kajian \citep{Patterson2025}. Penelitian ini tidak menerapkannya
karena data yang tersedia hanya berupa \textit{presence-only}. Pendekatan
ketiga adalah pembobotan sampel (\textit{instance weighting}). Pembobotan
sampel memberi bobot pada setiap sampel sumber agar sebaran sumber terbobot
menyerupai sebaran target, secara teoretis sebesar rasio kerapatan target
terhadap kerapatan sumber \citep{Zhuang2021,Kouw2021}.\ \citet{Saigusa2026}
menerapkan pembobotan sampel dengan rasio kerapatan yang diduga melalui uLSIF
pada transfer model antarwilayah. Rata-rata AUC uji pada 78 spesies hanya
berubah dari 0,633 menjadi 0,635. Namun, peningkatannya lebih tinggi saat
pergeseran sebaran besar dan catatan kemunculan sedikit \citep{Saigusa2026}.
Para penulis tersebut menyatakan hanya menggunakan uLSIF dan menyebut
\textit{Kernel Mean Matching} (KMM) sebagai alternatif \citep{Saigusa2026}.
KMM dipilih karena menduga bobot dengan menyamakan rata-rata sampel sumber dan
target pada ruang fitur kernel, tanpa menduga masing-masing kerapatan secara
langsung \citep{Zhuang2021}. Dalam pendugaan galat spasial saat terjadi
pergeseran sebaran, bobot rasio kerapatan dari estimasi kerapatan kernel
menjadi tidak stabil ketika sampel mengelompok. Pada kondisi tersebut, KMM
menghasilkan galat yang lebih rendah \citep{Serov2026}. Meskipun demikian,
pembobotan sampel mengasumsikan hubungan spesies dengan lingkungan (distribusi
bersyarat) sama pada kedua domain dan kondisi target berada dalam rentang
sumber \citep{Kouw2021}. Pembobotan sampel juga terutama sesuai apabila
perbedaan antara domain sumber dan target tidak terlalu besar
\citep{Zhang2024}. Pada tolok ukur metode adaptasi domain, pembobotan sampel
memberi peningkatan yang terbatas, bekerja paling baik pada model
teregularisasi, dan jarang memperburuk hasil (tahan terhadap transfer negatif)
\citep{Lalou2025}. Selain meminjam data dari wilayah lain, penelitian juga
dapat memperluas data kemunculan Indonesia sendiri dari sumber lain.

% DRAF P6 (Gemini Neon, 2026-11-21): sunting ke gaya sendiri.
Media sosial telah memasok data kemunculan berbagai spesies. Di Teluk Honda,
Palawan, Filipina, \citet{Araujo2019} menambah catatan individu hiu paus
melalui identifikasi foto dari operator lokal dan penambangan data unggahan
publik di Facebook, Instagram, Twitter, dan YouTube. Masyarakat menyumbang 36\%
dari 183 individu hiu paus yang teridentifikasi. Di Yunani, \citet{Taklis2020}
menelusuri Facebook, Twitter, dan YouTube pada 2017--2019 dan menemukan 116
unggahan tentang hiu. Mereka mencatat 22 kasus salah identifikasi oleh nelayan.
Temuan tersebut menunjukkan perlunya verifikasi data media sosial. Di Fiji,
sedikitnya dua peneliti meninjau setiap unggahan Facebook untuk memeriksa
identitas spesies, mutu media, tanggal, dan lokasi sebelum memakainya
\citep{Dehm2025}. Di Victoria, Australia, \citet{Beddoe2026} mencatat 3.209
laporan mamalia laut dari grup Facebook dan memverifikasi 1.124 laporan hingga
tingkat spesies setelah peninjauan pakar. Pada lima takson
(\textit{Tursiops} spp., orca, paus bungkuk, paus biru, dan paus sikat
selatan), Facebook menghasilkan proporsi sel kisi unik secara geografis yang
lebih tinggi daripada GBIF \citep{Beddoe2026}. Untuk delapan spesies lebah
Australia, \citet{Moore2024} menemukan bahwa penambahan catatan media sosial
(Flickr) atau sains komunitas (iNaturalist) pada model relung MaxEnt mengubah
sebaran prediksi dan bobot peubah lingkungan secara signifikan. Perubahan ini
lebih bergantung pada lokasi catatan tambahan daripada jumlahnya. Karena tidak
memiliki data uji independen, kajian ini hanya menunjukkan perubahan, bukan
apakah perubahan itu merupakan perbaikan. Kajian-kajian tersebut mengukur
jumlah titik, sel kisi unik, luas sebaran, atau perubahan prediksi. Namun,
kajian-kajian tersebut tidak mengukur sejauh mana catatan tambahan memperluas
cakupan lingkungan model kesesuaian habitat sambil mempertahankan kinerjanya.

% DRAF P7 (Gemini Neon, 2026-11-21): sunting ke gaya sendiri.
Berdasarkan uraian tersebut, penelitian ini membandingkan model kesesuaian
habitat hiu paus di perairan Indonesia yang hanya dilatih dengan data
Indonesia dengan skema yang memanfaatkan data Australia. Skema pemanfaatan data
Australia tersebut mencakup transfer langsung, adaptasi domain \textit{Kernel
Mean Matching} (KMM) berbasis pembobotan sampel, dan \textit{subsampling}
lingkungan. Penelitian ini juga menguji apakah penambahan data kemunculan dari
YouTube dan berita daring yang telah diverifikasi pakar memperluas cakupan
lingkungan model tanpa menurunkan kinerjanya.

\begin{comment}
Hiu paus (\textit{Rhincodon typus}) merupakan ikan terbesar di dunia dan
berstatus Terancam (\textit{Endangered}) dalam Daftar Merah IUCN
\citep{Pierce2025}. Di Indonesia, spesies ini dilindungi penuh sejak 2013
\citep{Herdianti2025}. Perairan Indonesia memiliki beberapa lokasi agregasi hiu
paus, antara lain Teluk Cenderawasih, Kaimana, Teluk Saleh, dan Teluk Tomini
\citep{Putra2026}. Di Bentang Laut Kepala Burung, hampir semua penampakan hiu
paus selama 2010--2023 tercatat di sekitar bagan \citep{Setyawan2025}.

Kajian habitat hiu paus di Indonesia umumnya berskala lokal dan berfokus pada
satu lokasi agregasi, misalnya Teluk Cenderawasih
\citep{Ranintyari2018,Manuhutu2021} dan Teluk Saleh \citep{Herdianti2025}.
Kajian tersebut menunjukkan bahwa suhu permukaan laut, klorofil-a, dan
kedalaman berkaitan dengan kemunculan hiu paus. Kajian berskala kepulauan yang
ada menggunakan data telemetri satelit \citep{Putra2026}. Data kemunculan pada
repositori terbuka seperti GBIF dan OBIS lebih mudah diakses, tetapi sebarannya
cenderung mengelompok di lokasi yang sering dikunjungi pengamat. Bias
pengambilan sampel seperti ini dapat menurunkan kemampuan model sebaran spesies
untuk diterapkan di luar wilayah pelatihannya \citep{Pili2025}.

Perairan Australia, terutama Ningaloo, memiliki data hiu paus yang lebih rapat
dari pemantauan jangka panjang dan telemetri satelit. Model habitat berbasis
telemetri dari Ningaloo memprediksi habitat yang sesuai hingga Laut Timor dan
perairan Indonesia \citep{Reynolds2017}. Namun, model yang dilatih di satu
wilayah belum tentu berkinerja baik ketika diterapkan di wilayah lain
\citep{Sahri2023}, terutama bila wilayah target memiliki kondisi lingkungan
yang tidak terwakili dalam data pelatihan \citep{Elith2010}. Adaptasi domain
dengan pembobotan sampel \textit{Kernel Mean Matching} (KMM) menyesuaikan bobot
data sumber agar sebarannya mendekati sebaran data target. Pada data spasial,
termasuk data kemunculan tumbuhan, KMM menghasilkan galat pendugaan yang lebih
rendah daripada model tanpa pembobotan \citep{Serov2026}.

Sumber lain yang dapat mengisi kekosongan data adalah unggahan media sosial dan
berita daring. Unggahan media sosial telah digunakan untuk mengidentifikasi
individu hiu paus di Filipina \citep{Araujo2019} dan mencatat kemunculan hiu di
Yunani \citep{Taklis2020}. Data ini memiliki keterbatasan karena spesies,
lokasi, dan tanggal tidak selalu dapat dipastikan, sehingga perlu diverifikasi
\citep{Dehm2025}. Penggabungan data oportunistik dengan data kemunculan lain
dapat memperbaiki pendugaan sebaran spesies yang datanya jarang dan mengelompok
\citep{Gonnerman2026}. Perubahan cakupan kondisi lingkungan akibat penambahan
data dapat dinilai dengan MESS \citep{Elith2010}.

Berdasarkan hal tersebut, penelitian ini memodelkan kesesuaian habitat hiu paus
di perairan Indonesia dengan memanfaatkan data Australia melalui adaptasi domain
KMM serta data media sosial yang diverifikasi pakar.

Revisi:
1. General ke spesifik, Australia dulu baru ke lokal
2. HSM tapi di luar negeri, Indonesia, dan negara tetangga. Apa bedanya penelitian kita dan penelitian mereka?
3. Sebelum rumusan masalah, latar belakang.
4. Kita harus tau masalah diselesaikan pake metode apa aja, kita juga harus memetakan juga dari paper paper yang relevan (see other methods?).
    - Kenapa kita pakai yang itu?
    - Bandingkan dengan metode selanjutnya
\end{comment}

\section{Rumusan Masalah}

Berdasarkan latar belakang tersebut, rumusan masalah dalam penelitian ini adalah
sebagai berikut:

\begin{enumerate}
    % Revisi 2026-11-21 (arah B, disetujui penulis): RM1 berfokus pada pengaruh
    % data Australia dengan empat skema; RM2 memuat syarat kinerja tidak turun.
    \item Bagaimana pengaruh pemanfaatan data kemunculan hiu paus
    (\textit{Rhincodon typus}) dari perairan Australia terhadap kinerja model
    kesesuaian habitat di perairan Indonesia apabila data tersebut digunakan
    melalui transfer langsung, adaptasi domain \textit{Kernel Mean Matching}
    (KMM), dan \textit{subsampling} lingkungan, dibandingkan dengan model yang
    hanya dilatih dengan data Indonesia?
    \item Sejauh mana penambahan data kemunculan dari media sosial dan berita
    daring yang telah diverifikasi pakar memperluas cakupan lingkungan model
    kesesuaian habitat hiu paus di perairan Indonesia tanpa menurunkan kinerja
    model?
    % Keputusan 2026-11-14: RM3 (pola musiman) dikeluarkan dari ruang lingkup.
    % Musim (NWM/TR1/SEM/TR2) tetap dipakai sebagai irisan data pada metode,
    % tetapi tidak menjadi tujuan penelitian dan tidak boleh dijanjikan di Tujuan.
    % \item Bagaimana pola sebaran kesesuaian habitat musiman hiu paus
    % (\textit{Rhincodon typus}) pada empat musim muson di perairan Indonesia
    % (\textit{Northwest Monsoon}/NWM, \textit{Transition} 1/TR1,
    % \textit{Southeast Monsoon}/SEM, dan \textit{Transition} 2/TR2) serta
    % peubah oseanografi dan batimetri apa yang paling menentukan kesesuaian
    % habitat tersebut?
\end{enumerate}

\begin{comment}
Rumusan masalah merupakan pernyataan ringkas mengenai masalah yang akan
diselesaikan dan cara mengatasinya untuk menjawab tujuan tugas akhir. Masalah
yang diteliti dapat dirumuskan karena berbagai sebab, seperti adanya
kesenjangan (\textit{gap}), tantangan, kesangsian, ketidakjelasan, dan
keingintahuan secara akademik yang berkaitan dengan fenomena alam, sosial, dan
ekonomi. Dalam pernyataan ringkas tersebut harus tercakup pendekatan yang
digunakan dalam perumusan masalah.
\end{comment}


\section{Tujuan}

Berdasarkan rumusan masalah yang telah diuraikan, penelitian kesesuaian habitat
hiu paus (\textit{Rhincodon typus}) ini bertujuan untuk:

\begin{enumerate}
    \item Membandingkan kinerja model kesesuaian habitat hiu paus di perairan
    Indonesia pada empat skema pelatihan, yaitu model lokal Indonesia, transfer
    langsung dari Australia, adaptasi domain KMM, dan \textit{subsampling}
    lingkungan, menggunakan algoritme MaxEnt, XGBoost, dan \textit{Random
    Forest}. Kinerja diukur dengan indeks Boyce kontinu sebagai ukuran utama
    serta AUC-ROC dan TSS sebagai ukuran pelengkap, melalui validasi silang
    blok spasial lima lipatan (\textit{spatial block cross-validation}).
    \item Mengukur perubahan luas \textit{area of applicability} (AOA), luas
    wilayah ekstrapolasi MESS (\textit{Multivariate Environmental Similarity
    Surface}), dan indeks Boyce kontinu model di perairan Indonesia akibat
    penambahan data media sosial terverifikasi pada data GBIF dan OBIS.
\end{enumerate}

\begin{comment}
Pernyataan tujuan tugas akhir ialah pernyataan singkat dan jelas tentang tujuan
yang akan dicapai sebagai upaya pemecahan masalah maupun memahami gejala
(fenomena) yang dijelaskan dalam latar belakang. Tujuan merupakan pemandu atau
arah untuk merencanakan dan melaksanakan kajian tugas akhir. Tujuan penelitian
ditulis dengan memilih kata kerja yang hasilnya dapat diukur dan dilihat,
seperti: menguraikan, menerangkan, membuktikan, menjajaki, menguji,
membuktikan, atau menerapkan suatu gejala, konsep atau dugaan, atau bahkan
membuat suatu prototipe. Jangan menggunakan kata kerja mengetahui atau
memahami. Masalah dan tujuan penelitian harus terkait dan konsisten.
\end{comment}

\section{Manfaat}

Manfaat merupakan dampak positif (kegunaan) dari hasil karya ilmiah tugas akhir
bagi bidang ipteks, pembangunan, dan masyarakat. Manfaat utama dari hasil tugas
akhir adalah menambah khasanah ilmu pengetahuan dalam bentuk pustaka sebagai
sumber acuan/referensi untuk pengembangan ipteks, para pengambil keputusan baik
di industri maupun pemerintah dan lembaga untuk menyusun kebijakan baru, serta
masyarakat umum. Manfaat dinyatakan dengan kata kerja yang lugas dan logis.

\section{Ruang Lingkup}

Penelitian ini dibatasi pada hal-hal berikut.

\begin{enumerate}
    \item Penelitian memodelkan kesesuaian habitat satu spesies, yaitu hiu paus
    (\textit{Rhincodon typus}), di perairan Indonesia sebagai domain target.
    Perairan Australia digunakan sebagai domain sumber untuk adaptasi domain.
    \item Data kemunculan berasal dari GBIF dan OBIS pada periode 2010--2026,
    dengan ketidakpastian koordinat paling besar 10 km, serta data media sosial
    yang telah diverifikasi pakar. Hanya unggahan berstatus valid yang
    digunakan. Data media sosial dan berita daring dikumpulkan dari YouTube dan
    laporan berita daring. Instagram, TikTok, dan X tidak dicakup karena
    pengumpulan otomatis tidak sesuai dengan ketentuan platform. Data dari
    iNaturalist dan Sharkbook tidak digunakan dalam penelitian ini.
    % Opsi 2 (belum dipilih, keputusan 2026-11-16): ganti dua kalimat tentang
    % Instagram/TikTok/X di atas dengan: "Instagram, TikTok, dan X dicakup
    % melalui ekspor manual atau API resmi, lalu diverifikasi pakar."
    % Syarat: ada minimal satu baris valid per platform; jika tidak, platform itu
    % tidak boleh disebut. API resmi: X berbayar (sekitar 0,005 USD per unggahan),
    % TikTok Research API tidak tersedia untuk peneliti Indonesia, Instagram
    % memerlukan verifikasi bisnis dan App Review. Bantuan LLM hanya untuk
    % penyaringan awal dan ekstraksi kolom; verifikasi tetap oleh pakar.
    \item Peubah lingkungan berupa suhu permukaan laut, tinggi muka laut,
    klorofil-a, produksi primer bersih, kedalaman, kemiringan dasar laut, dan
    jarak ke garis pantai serta isobat 200, 300, dan 500 m, pada resolusi
    1/12$^{\circ}$ (sekitar 9,2 km). Suhu dasar laut tidak digunakan karena
    mencerminkan massa air dalam, bukan habitat makan di lapisan permukaan.
    \item Model yang dibandingkan adalah MaxEnt, XGBoost, dan
    \textit{Random Forest} pada empat skema pelatihan (lokal Indonesia,
    transfer langsung, KMM, dan \textit{subsampling} lingkungan), dengan
    validasi silang blok spasial lima lipatan.
    \item Musim muson hanya digunakan sebagai pengelompokan data dan peubah
    dinamis. Penelitian tidak menganalisis pola sebaran musiman.
    \item Cakupan lingkungan diukur terutama dengan AOA
    \citep{MeyerPebesma2021}. MESS menjadi ukuran pendukung karena hanya
    mendeteksi nilai di luar rentang satu peubah dan tidak mendeteksi
    perubahan korelasi atau kombinasi baru antarpeubah \citep{Elith2010}.
\end{enumerate}

\begin{comment}
Tugas akhir sering kali dihadapkan pada keterbatasan data, dana, waktu, metode,
bahkan teori. Oleh karena itu, tugas akhir perlu dengan tegas menunjukkan ruang
lingkup dengan mempertimbangkan keterbatasan tersebut.
\end{comment}

\section{Hipotesis}

\begin{comment}
Hipotesis dapat ditulis secara eksplisit atau tersirat sesuai bidang ipteks yang
relevan. Hipotesis dapat menjadi bagian dari pendahuluan (untuk bidang sains,
teknik, dan kesehatan), atau merupakan bagian akhir dari tinjauan pustaka
(untuk bidang sosial, ekonomi, dan humaniora).
\end{comment}

Hipotesis penelitian ini adalah sebagai berikut:

\begin{enumerate}
    \item Pemanfaatan data Australia meningkatkan indeks Boyce kontinu model di
    perairan Indonesia dibandingkan model lokal, dan peningkatan dari KMM lebih
    besar pada MaxEnt daripada pada XGBoost dan \textit{Random Forest}
    \citep{Lalou2025}.
    \item Penambahan data media sosial terverifikasi memperluas AOA dan
    mempersempit wilayah ekstrapolasi MESS tanpa menurunkan indeks Boyce
    kontinu.
\end{enumerate}
